% ===========================================================================
% mappe/mechanik-mappe.tex — Darstellungsebene der Charaktermappe
%
% Dieselbe Trennung wie beim Heldenbogen: hier steht nur, WIE etwas aussieht.
% WAS drinsteht, kommt aus helden/<name>.tex, und WO die Blätter landen,
% entscheidet mappe.tex.
%
% Der Kern ist die Umgebung mappenblatt: sie baut EIN A4-Feld von genau
% 210 x 297 mm als Box. Die digitale Fassung setzt eines pro Seite, die
% Druckfassung zwei davon nebeneinander auf A3 quer. Deshalb darf hier
% nichts vom Seitenrand abhängen — alles wird absolut in das Feld gesetzt.
%
% Koordinaten in Millimeter von der linken OBEREN Feldecke, wie im ganzen
% Projekt. Erreicht wird das mit y=-1mm; das dreht die y-Achse um und hat
% zwei Folgen, die weiter unten je an ihrer Stelle stehen: rotate=90 und
% rotate=-90 tauschen den Sinn, und yshift zählt weiter nach oben.
% ===========================================================================

\makeatletter

% --- Werte einsammeln ------------------------------------------------------
% Gleiche Schnittstelle wie mechanik.tex, damit helden/<name>.tex
% unverändert durch beide Dokumente läuft.
\newcommand{\wert}[2]{\expandafter\gdef\csname wert@#1\endcsname{#2}}
\newcommand{\wertvon}[1]{\ifcsname wert@#1\endcsname\csname wert@#1\endcsname\fi}
\newcommand{\wertdaP}[1]{%
  \ifcsname wert@#1\endcsname\expandafter\@firstoftwo\else\expandafter\@secondoftwo\fi}
% Erster gesetzter Wert aus zwei Namen, sonst leer. Damit greift die Mappe
% auf die Felder des Heldenbogens zurück: mappe_name, sonst s1_name.
\newcommand{\wertoder}[2]{\wertdaP{#1}{\wertvon{#1}}{\wertvon{#2}}}
% Wert oder Vorgabe.
\newcommand{\mappewahl}[2]{\wertdaP{#1}{\wertvon{#1}}{#2}}

% Die Heldendateien dürfen die Bogenauswahl setzen. Für die Mappe ist das
% ohne Bedeutung — hier verschluckt, damit \input nicht scheitert.
\newcommand{\bogennur}[1]{}
\newcommand{\bogenweglassen}[1]{}

% --- Farben ----------------------------------------------------------------
% Am HvS-Charakterbogen gemessen (300 dpi, Querschnitt durch den Rahmen).
\definecolor{bandcreme}{HTML}{D5C7C4}
\definecolor{bandkontur}{HTML}{412A1E}
\definecolor{schrift}{HTML}{241A14}
\definecolor{pergament}{HTML}{FCFAF6}

% Die Akzentfarbe kommt aus der Heldendatei: \wert{mappe_akzentfarbe}{AD6ED0}.
% Fehlt sie, bleibt der Rahmen im gedeckten Braun des Bandes — nicht
% farblos, nur ohne Akzent.
\newcommand{\akzentsetzen}{%
  \wertdaP{mappe_akzentfarbe}{%
    % \definecolor expandiert seine Angabe nicht selbst — erst den Aufruf
    % zusammenbauen, dann ausführen.
    \edef\akzent@ruf{\noexpand\definecolor{akzent}{HTML}{\wertvon{mappe_akzentfarbe}}}%
    \akzent@ruf
  }{\colorlet{akzent}{bandkontur!35!bandcreme}}}

% --- Flechtband ------------------------------------------------------------
% Am Vorbild eingemessen. Maße in mm von der Blattkante:
%   1,75--2,10  dunkle Kontur   \
%   2,15--2,85  Creme-Körper     > Strang A  (Mitte 2,45  Dicke 1,40)
%   2,95--3,15  dunkle Kontur   /
%   3,20--4,35  AKZENTFARBE       (die Lücke zwischen den Strängen!)
%   4,40--4,60  dunkle Kontur   \
%   4,65--5,25  Creme-Körper     > Strang B  (Mitte 5,00)
%   5,35--5,60  dunkle Kontur   /
%   6,85--8,35  glatte Innenlinie (Creme 7,20--7,95)
% Kettenmaß 18,12 mm, davon 11,00 mm gerader Balken.
%
% Der Witz des Motivs: es sind ZWEI parallele Bänder. Die Akzentfarbe ist
% nichts als der Zwischenraum. Auf der Geraden wird daraus ein langer
% Balken, zwischen zwei Kreuzungen eine spitze Linse. Ein Farbwert genügt,
% um den ganzen Rahmen umzustimmen.
\def\bandmitte{3.68}
\def\strangversatz{1.23}
\def\strangbauch{1.78}
\def\strangdick{1.40}
\def\strangcreme{0.70}
\def\innenlinie{7.60}
\def\innendick{1.50}
\def\innencreme{0.75}
\def\periode{18.12}
\def\balken{11.00}

% Strangkurve. \pp \bb \kk \np müssen vorher gesetzt sein — sie DÜRFEN nicht
% hier gesetzt werden: eine Zuweisung mitten im Pfad bricht TikZ ab
% ("Giving up on this path. Did you forget a semicolon?").
% #1 = Vorzeichen des Versatzes.
\newcommand{\strangvor}[1]{%
  (0,{#1*\strangversatz})%
  \foreach \i in {0,...,\the\numexpr\np-1\relax} {%
    -- ({\i*\pp+\bb},{#1*\strangversatz})
    .. controls ({\i*\pp+\bb+0.20*\kk},{#1*\strangversatz})
            and ({\i*\pp+\bb+0.27*\kk},{-#1*\strangbauch})
     .. ({\i*\pp+\bb+0.50*\kk},{-#1*\strangbauch})
    .. controls ({\i*\pp+\bb+0.73*\kk},{-#1*\strangbauch})
            and ({\i*\pp+\bb+0.80*\kk},{#1*\strangversatz})
     .. ({\i*\pp+\pp},{#1*\strangversatz})%
  }%
}
% Derselbe Weg rückwärts. Gebraucht, um die Akzentfläche als geschlossenen
% Pfad zu bekommen: Strang A hin, Strang B zurück.
\newcommand{\strangrueck}[1]{%
  \foreach \j in {1,...,\np} {%
    -- ({(\np-\j)*\pp+\pp},{#1*\strangversatz})
    .. controls ({(\np-\j)*\pp+\bb+0.80*\kk},{#1*\strangversatz})
            and ({(\np-\j)*\pp+\bb+0.73*\kk},{-#1*\strangbauch})
     .. ({(\np-\j)*\pp+\bb+0.50*\kk},{-#1*\strangbauch})
    .. controls ({(\np-\j)*\pp+\bb+0.27*\kk},{-#1*\strangbauch})
            and ({(\np-\j)*\pp+\bb+0.20*\kk},{#1*\strangversatz})
     .. ({(\np-\j)*\pp+\bb},{#1*\strangversatz})
    -- ({(\np-\j)*\pp},{#1*\strangversatz})%
  }%
}

% Eine Kante: läuft von (0,0) nach (#1,0), +y zeigt nach außen.
% #2 = Anzahl Kettenglieder. Die Konturen ALLER Stränge zuerst, dann alle
% Creme-Körper — sonst deckt die Kontur des zweiten Strangs den ersten ab.
\newcommand{\flechtkante}[2]{%
  \def\np{#2}%
  \pgfmathsetmacro{\pp}{#1/#2}%
  \pgfmathsetmacro{\bb}{\pp*\balken/\periode}%
  \pgfmathsetmacro{\kk}{\pp-\bb}%
  \fill[akzent] \strangvor{-1} \strangrueck{1} -- cycle;
  \foreach \sg in {1,-1} {%
    \draw[bandkontur,line width=\strangdick mm,line cap=round] \strangvor{\sg};}%
  \foreach \sg in {1,-1} {%
    \draw[bandcreme,line width=\strangcreme mm,line cap=round] \strangvor{\sg};}%
}

% Ecke: beide Stränge gehrt um die Ecke, dazwischen die Akzentfläche.
% Scheitel in (0,0), die Kanten laufen nach +x und +y. #1 = Schenkellänge.
\newcommand{\flechtecke}[1]{%
  \fill[akzent] (#1,{-\strangversatz}) -- ({-\strangversatz},{-\strangversatz})
             -- ({-\strangversatz},#1) -- ({\strangversatz},#1)
             -- ({\strangversatz},{\strangversatz}) -- (#1,{\strangversatz}) -- cycle;
  \foreach \s in {\strangversatz,-\strangversatz} {%
    \draw[bandkontur,line width=\strangdick mm,line join=round,line cap=round]
      (#1,{-\s}) -- ({-\s},{-\s}) -- ({-\s},#1);}%
  \foreach \s in {\strangversatz,-\strangversatz} {%
    \draw[bandcreme,line width=\strangcreme mm,line join=round,line cap=round]
      (#1,{-\s}) -- ({-\s},{-\s}) -- ({-\s},#1);}%
}

% Der ganze Rahmen. Wird in einem Feld mit y nach UNTEN aufgerufen.
% Der Einzug ist so gewählt, dass die Kette waagerecht in 10 und senkrecht
% in 15 gleiche Glieder aufgeht — 18,3 bzw. 17,96 mm gegen die gemessenen
% 18,12. Das ist der einzige Punkt, an dem vom Maß abgewichen wird, und der
% Grund dafür ist zwingend: ein halbes Glied in der Ecke fällt auf.
\def\kanteneinzug{13.6}
\newcommand{\flechtrahmen}{%
  \pgfmathsetmacro{\ez}{\kanteneinzug}
  \pgfmathsetmacro{\lw}{210-2*\ez}
  \pgfmathsetmacro{\lh}{297-2*\ez}
  \pgfmathsetmacro{\el}{\ez-\bandmitte}
  % glatte Innenlinie als EIN umlaufendes Rechteck — so schließen die Ecken
  \draw[bandkontur,line width=\innendick mm]
    (\innenlinie,\innenlinie) rectangle ({210-\innenlinie},{297-\innenlinie});
  \draw[bandcreme,line width=\innencreme mm]
    (\innenlinie,\innenlinie) rectangle ({210-\innenlinie},{297-\innenlinie});
  % Kanten. Das +y der Kante zeigt jeweils nach außen. Weil die y-Achse
  % dieses Feldes nach unten läuft, ist der Drehsinn getauscht: rotate=90
  % dreht das +x der Kante nach OBEN.
  \begin{scope}[shift={(\ez,\bandmitte)},yscale=-1]         \flechtkante{\lw}{10}\end{scope}
  \begin{scope}[shift={(210-\ez,297-\bandmitte)},xscale=-1] \flechtkante{\lw}{10}\end{scope}
  \begin{scope}[shift={(\bandmitte,297-\ez)},rotate=90]     \flechtkante{\lh}{15}\end{scope}
  \begin{scope}[shift={(210-\bandmitte,\ez)},rotate=-90]    \flechtkante{\lh}{15}\end{scope}
  % Ecken
  \begin{scope}[shift={(\bandmitte,\bandmitte)}]                             \flechtecke{\el}\end{scope}
  \begin{scope}[shift={(210-\bandmitte,\bandmitte)},xscale=-1]               \flechtecke{\el}\end{scope}
  \begin{scope}[shift={(\bandmitte,297-\bandmitte)},yscale=-1]               \flechtecke{\el}\end{scope}
  \begin{scope}[shift={(210-\bandmitte,297-\bandmitte)},xscale=-1,yscale=-1] \flechtecke{\el}\end{scope}
}

% Kopfzier: dort, wo beim Vorbild das Verlagslogo steht. Das ist eine
% eingetragene Marke und wird NICHT nachgezeichnet. Stattdessen ein
% freistehendes Stück derselben Kette — es nimmt die Akzentfarbe mit, was
% das Logo nie getan hätte.
% #1 = Mitte x, #2 = Mitte y, #3 = Breite, #4 = Glieder
\newcommand{\kopfzier}[4]{%
  \begin{scope}[shift={({#1-#3/2},#2)}]
    \flechtkante{#3}{#4}
  \end{scope}}

% --- Pergament -------------------------------------------------------------
% Fläche und Fußkasten sind aus dem Scriptorium-Baukasten abgeleitet, siehe
% werkzeuge/pergament.py. Sie liegen in grafiken/ wie jede andere
% Baukastengrafik und sind Verlagsmaterial — sie werden nicht weitergegeben,
% wohl aber das fertige PDF (siehe RECHTLICHES auf dem linken Innenblatt).
%
% Der Pfad ist relativ zum Arbeitsverzeichnis, und das ist bogen/ — die
% Bauskripte setzen es dorthin. Deshalb ../grafiken/, nicht grafiken/.
\newcommand{\pergamentgrund}{%
  \node[anchor=north west,inner sep=0pt] at (0,0)
    {\includegraphics[width=210mm,height=297mm]{../grafiken/mappe-pergament-a4.jpg}};}

% Ein Pergamentkasten mit gerissenem Rand, wie am Fuß der Vorbildseiten.
% #1 x, #2 y (linke obere Ecke, mm), #3 Breite in mm, #4 Inhalt.
%
% Die HÖHE wird nicht angegeben, sondern am Inhalt gemessen. Ein fester
% Wert lässt bei kurzem Text ein leeres Stück Pergament stehen, und das
% sieht nach Versehen aus.
% Der Einzug ist waagerecht groesser als senkrecht. Grund: der gerissene
% Rand der Vorlage wird beim Einpassen mitskaliert. Die Vorlage ist
% 1010 x 599 Pixel, der Fusskasten aber 182 x 28 mm -- waagerecht wird der
% Rand dadurch auf gut 4 mm gestreckt, senkrecht auf gut einen. Mit einem
% einheitlichen Einzug steht der Text seitlich in der Faserung.
\newlength{\kasteneinzug}
\setlength{\kasteneinzug}{12mm}
\newlength{\kasteneinzughoch}
\setlength{\kasteneinzughoch}{6mm}
\newsavebox{\kastenbox}
\newcommand{\pergamentkasten}[4]{%
  % pgfinterruptpicture ist hier PFLICHT. Ein \sbox mit einer minipage
  % darin, direkt in einem tikzpicture aufgerufen, kommt LEER zurueck --
  % gemessen: ht 2,5 pt, dp 0 pt bei richtiger Breite, ohne jede
  % Fehlermeldung. Der Kasten schrumpft dann auf die doppelte Randbreite
  % und der Text fehlt einfach. pgfinterruptpicture stellt fuer den
  % Augenblick den gewoehnlichen LaTeX-Satzzustand her; \global, weil die
  % Umgebung eine Gruppe ist und die Box sie sonst nicht ueberlebt.
  \begin{pgfinterruptpicture}%
    \global\setbox\kastenbox\hbox{%
      \begin{minipage}{\dimexpr#3mm-2\kasteneinzug\relax}%
        \raggedright\color{schrift}#4\par
      \end{minipage}}%
  \end{pgfinterruptpicture}%
  \node[anchor=north west,inner sep=0pt] at (#1,#2)
    {\includegraphics[width=#3mm,
       height=\dimexpr\ht\kastenbox+\dp\kastenbox+2\kasteneinzughoch\relax]
       {../grafiken/mappe-pergament-kasten.png}};
  % yshift zählt in Leinwandkoordinaten, dort ist +y oben — daher das Minus.
  \node[anchor=north west,inner sep=0pt]
    at ([xshift=\kasteneinzug,yshift=-\kasteneinzughoch]#1,#2) {\usebox\kastenbox};}

% --- Schrift ---------------------------------------------------------------
% Cinzel für die Auszeichnung, EB Garamond für den Lesetext. Beide liegen in
% TeX Live und laufen unter pdflatex.
%
% Zwei Dinge, die daran nicht zu ändern sind:
% - Die Titelschrift des Vorbilds ist ein Rasterbild. Sie ist damit nicht
%   greifbar, auch nicht über die Schriftliste des PDF — dort steht nur
%   Gentium Basic, und das nur für den Fußtext.
% - Der Baukasten schreibt Gentium Basic als Fließtext vor. Ein
%   pdflatex-fähiges Gentium gibt es in TeX Live nicht mehr
%   ("tlmgr: package gentium-tug not present in repository"). EB Garamond
%   ist der nächste Verwandte, der da ist.
\newcommand{\auszeichnung}{\cinzelblack}
\newcommand{\lesetext}{\rmfamily}

% ACHTUNG, zweimal eingetreten:
%
% 1. Die Schriftwahl muss als font=-Option des Knotens stehen, NICHT im
%    Knotentext. Steht sie im Text, gilt sie nur für die erste Zeile: nach
%    einem \\ in einem align-Knoten fällt die Größe auf die Grundschrift
%    zurück, ohne Fehlermeldung. Der Titel sieht dann aus wie ein Name in
%    40 pt mit einem Nachnamen in 10 pt.
% 2. Mehrabsätziger Text braucht ein \par am Ende. Der Zeilenabstand eines
%    Absatzes wird bei seinem \par festgelegt; fehlt es, setzt TeX den
%    letzten Absatz mit dem Zeilenabstand von AUSSERHALB des Knotens. In
%    einem Text aus drei Absätzen sind dann zwei locker und einer eng.
%    \pergamentkasten und \mappetextblock hängen das \par selbst an.

% Zeile in Akzentfarbe mit dunkler Kontur, wie der Name auf dem Vorbild.
% Umgesetzt als Versatzkranz: derselbe Text achtmal in der Konturfarbe
% ringsum, dann einmal in der Akzentfarbe darüber. Umständlicher als ein
% echtes Umranden, dafür ohne \pdfliteral im Textmodus — und bei 40 pt ist
% das Ergebnis nicht zu unterscheiden.
% #1 x, #2 y (Mitte oben), #3 Konturradius in mm, #4 Schriftwahl, #5 Inhalt
\newcommand{\konturzeile}[5]{%
  \foreach \dx/\dy in {-1/0,1/0,0/-1,0/1,-0.71/-0.71,0.71/-0.71,-0.71/0.71,0.71/0.71} {%
    \node[anchor=north,align=center,font=#4,text=bandkontur]
      at ({#1+\dx*#3},{#2+\dy*#3}) {#5};}%
  \node[anchor=north,align=center,font=#4,text=akzent] at (#1,#2) {#5};}

% Lesetext-Block. #1 x, #2 y (linke obere Ecke), #3 Breite mm,
% #4 Schriftwahl, #5 Ausrichtung (justify/center), #6 Inhalt.
\newcommand{\mappetextblock}[6]{%
  \node[anchor=north west,text=schrift,text width=#3mm,align=#5,font=#4]
    at (#1,#2) {#6\par};}

% --- Stellschrauben mit Vorgabewert ----------------------------------------
% Alle über die Heldendatei übersteuerbar, alle mit brauchbarer Vorgabe.
% Zweck: ein langer Name braucht eine kleinere Schrift, und das soll nicht
% dazu zwingen, das Layout anzufassen.
\newcommand{\mappenamengroesse}{\mappewahl{mappe_namengroesse}{40}}
\newcommand{\mappenamenzeile}{\mappewahl{mappe_namenzeile}{45}}
\newcommand{\mappeprofgroesse}{\mappewahl{mappe_professionsgroesse}{21}}
\newcommand{\mappehintergrundtitel}{\mappewahl{mappe_hintergrundtitel}{HINTERGRUND}}
\newcommand{\mappecopyright}{\mappewahl{mappe_copyright}{[Jahr] [Name]}}

% --- Heldenbild ------------------------------------------------------------
% #1 bevorzugter Wertename, #2 Rückfall, dann x, y, Breite, Höhe (mm).
%
% Das Bild sitzt auf der UNTERKANTE des Kastens, nicht mittig darin. Eine
% stehende Figur, die in der Luft hängt, fällt sofort auf; ein Bild, das
% oben Luft lässt, nicht. Fehlt Wert oder Datei, bleibt die Fläche leer und
% der Lauf schreibt nur einen Hinweis ins Log — wie bei jedem Feld.
\newcommand{\mappebild}[6]{%
  \edef\mappe@bild{\wertoder{#1}{#2}}%
  \ifx\mappe@bild\empty
    \message{^^JHINWEIS: kein Bild gesetzt (#1 / #2), Bildflaeche bleibt leer.^^J}%
  \else
    \IfFileExists{\mappe@bild}{%
      \node[anchor=south,inner sep=0pt] at ({#3+#5/2},{#4+#6})
        {\includegraphics[width=#5mm,height=#6mm,keepaspectratio]{\mappe@bild}};%
    }{\message{^^JHINWEIS: Bilddatei nicht gefunden: \mappe@bild^^J}}%
  \fi}

% --- Ein Blatt -------------------------------------------------------------
% Baut ein A4-Feld von genau 210 x 297 mm als Box. y läuft nach unten,
% (0,0) ist oben links.
\newsavebox{\mappenbox}
\newenvironment{mappenblatt}{%
  \akzentsetzen
  \begin{lrbox}{\mappenbox}%
  \begin{tikzpicture}[x=1mm,y=-1mm,inner sep=0pt,outer sep=0pt]
    \useasboundingbox (0,0) rectangle (210,297);
    \begin{scope}
      \clip (0,0) rectangle (210,297);
      \pergamentgrund
}{%
    \end{scope}
  \end{tikzpicture}%
  \end{lrbox}%
  \usebox{\mappenbox}}

\makeatother
